\documentclass{article}
\usepackage[T1]{fontenc}
\usepackage{lmodern}
\usepackage{graphicx}
\usepackage{amsmath,amssymb}
\usepackage[a4paper,margin=2cm]{geometry}
\usepackage[absolute,overlay]{textpos}
\setlength{\TPHorizModule}{1mm}
\setlength{\TPVertModule}{1mm}
\usepackage[indonesian]{babel}
\usepackage{hyperref}
\setlength{\emergencystretch}{2em}
% unit: Halaman 7

\begin{document}

% === Halaman 7 (python/native) ===

• Perhatikan bahwa a = 2 bukan akar primitif dari 7 karena 

21 (mod 7) = 2;

22 (mod 7) = 4;

23 (mod 7) = 1

24 (mod 7) = 2;

25 (mod 7) = 4;

26 (mod 7) = 1

• Nilai-nilai yang dihasilkan dari 21, 22, …, 26 tidak semuanya berbeda dan tidak 

mencakup semua nilai di dalam modulus 7.

• Untuk menemukan semua akar primitif dari p, kita harus mencoba semua 

bilangan bulat dari 2, 3, …  

7

\end{document}
